\subsection{Functoriality, congruence descent, and reconstruction}

\begin{statement}{B.15}{Proved}{The exact descent criterion and the corrected reconstruction formalism}
We retain the definitions and topology of B.13, including the
recorded numerical coordinates. Let
\(q:C\twoheadrightarrow C'\) be a continuous epimorphism of
profinite groups carrying the ordered generators to
\((\rho',\sigma')\) and commuting with the involutions \(\tau\).
Use either \(D=B\) in both groups or \(D=Z\) in both groups, and put
\(N=\ker q\). The maximal subgroup on which quotient automorphisms
are defined is
\[
 D_q=\{d\in D:u_d(N)=N\}.
 \tag{B.15.1}
\]
It is closed, contains \(j(L_0)\) and \(t\), and has a canonical
continuous homomorphism
\(q_D:D_q\to D'\), preserving every numerical coordinate. Write
\(H=q_D(D_q)\), a closed subgroup of \(D'\).

The following conditions are equivalent: \(D_q=D\); every effective
source numerical solution maps to an effective target solution;
and the quotient automorphism \(q_D(d)\) exists for every \(d\in D\).
Replacing normalization of \(L_0\) by conjugacy of \(L_0\) does not
change the corresponding kernel-stability criterion. These are the
precise equivalent conditions of C3:570, (89)--(92), folio~588.
Uniqueness of the resulting map of the enlarged groups includes
compatibility with both specified embeddings \(C\to G_D\) and
\(D\to G_D\), as added in C3:571; compatibility with the induced
inner actions alone is not being substituted for this condition.

Assume, for the next formulas, that
\(L_0\cap Z(C)=L'_0\cap Z(C')=1\). Form the B.13 groups using
\(D_q\) and \(H\), respectively, and set
\[
 A_q=\ker(D_q\to H),\qquad L_N=L_0\cap N,
 \qquad K_G=\ker(G_{D_q}\to G_H).
\]
The maps of \(D,\gamma,G\) onto their indicated targets are
continuous epimorphisms, and their exact kernels are
\[
\begin{split}
 K_\gamma&:=\ker(\gamma_{D_q}\to\gamma_H)
                  \simeq A_q/j(L_N),\\
 K_G&\simeq (N\rtimes A_q)/
       \{(\ell^{-1},j(\ell)):\ell\in L_N\},\\
 1&\longrightarrow N\longrightarrow K_G
          \longrightarrow K_\gamma\longrightarrow1.
\end{split}
\tag{B.15.2}
\]
For \(x=\rho,\sigma\), the induced map of signed normalizers has
kernel
\[
 K_x:=\ker(N_x^*\to N_{x'}^*)=C_{K_G}(x).
 \tag{B.15.3}
\]
It need not be onto. If \((g',e)\in N_{x'}^*\) and \(g\) is any
lift of \(g'\), the exact obstruction is
\[
 (g',e)\text{ lifts}\quad\Longleftrightarrow\quad
 gxg^{-1}\text{ is }K_G\text{-conjugate to }x^e.
 \tag{B.15.4}
\]
Every nonempty fibre is a right torsor under \(C_{K_G}(x)\).
The maps of \(\Gamma,\mathcal M,S_0\Gamma,S_0\mathcal M\) are the
componentwise fibre-product maps. In particular their kernels are
\[
\begin{aligned}
 K_\Gamma&=K_\rho\times_{K_\gamma}K_\sigma,&
 K_{\mathcal M}&=K_\Gamma\times_{K_\gamma}K_G,\\
 K_{S_0\Gamma}&=A_q\times_{K_\gamma}K_\Gamma,&
 K_{S_0\mathcal M}&=A_q\times_{K_\gamma}K_{\mathcal M}.
\end{aligned}
\tag{B.15.5}
\]
No surjectivity of these last four maps follows merely from
surjectivity of \(q\).

The subgroup variant of C3:571--572, (94), has the following precise
form. For a closed \(D_0\subset D\) containing \(t\), set
\[
\begin{split}
 L(D_0)&=\{\ell\in L_0:j(\ell)\in D_0\},\qquad
 \gamma_0=D_0/j(L(D_0)),\\
 G_0&=CD_0\subset G_D,\qquad
 G_0\simeq(C\rtimes D_0)/\operatorname{anti}L(D_0),\\
 N_x^{0,*}&=N_x^*\cap(G_0\times\{\pm1\}).
\end{split}
\tag{B.15.6}
\]
Then \(1\to C\to G_0\to\gamma_0\to1\) is exact, and the
remaining groups are the B.13 fibre products over \(\gamma_0\).
The kernel of \(D_0\to\gamma_0\) is \(j(L(D_0))\), which equals
\(j(L_0)\) only if the latter is contained in \(D_0\).
For descent replace \(D_0\) by \(D_0\cap D_q\).
The distinguished element required here is \(t\in D\); the
manuscript's margin reading \(\tau'_Z\) cannot mean
\(t'=\sigma^{-1}t\), which need not belong to \(D\).
This construction keeps the full conjugating group \(C\). The
different restriction of the conjugating group in \S49.VI is
treated separately in B.17--B.18.

For the particular map
\[
 q:\widehat{\SL_2(\Z)}\twoheadrightarrow\SL_2(\widehat\Z),
 \quad
 \rho=\begin{pmatrix}0&1\\-1&1\end{pmatrix},\quad
 \sigma=\begin{pmatrix}0&-1\\1&0\end{pmatrix},
 \tag{B.15.7}
\]
we have \(D_q=D\), for either choice of \(D\). This is the earlier
\(\sigma\)-convention of C3:572--573 and B.13.
More explicitly, if \(q(u(\sigma))_{11}=a\) and the recorded
multiplier is \(\mu\), then
\[
 q\circ u=\operatorname{Int}(g(a,\mu))\circ q,
 \qquad g(a,\mu)=\begin{pmatrix}\mu&a\mu\\0&1\end{pmatrix}.
 \tag{B.15.8}
\]
The restriction \(L_0\to L'_0\) is an isomorphism, so here
\(L_N=1\), \(K_\gamma=A_q\), and \(K_G=N\rtimes A_q\).
Thus congruence-kernel stability is proved; it is not one of the
unresolved lifting assertions in B.8 or D.7.

Finally, the reconstruction at the end of \S49, C3:691--696,
folios~648--650, has a general form. Let \(P,T\) be profinite
groups, \(J\triangleleft T\) closed, \(\theta:T\to\Aut(P)\)
a continuous action, and \(\phi:J\to P\) a continuous homomorphism
satisfying
\[
 \theta(j)=\operatorname{Int}(\phi(j)),\qquad
 \phi(tjt^{-1})=\theta(t)(\phi(j)).
 \tag{B.15.9}
\]
Then
\[
 \mathcal E=(P\rtimes T)/
      \{(\phi(j)^{-1},j):j\in J\},\qquad
 1\longrightarrow P\longrightarrow\mathcal E
     \longrightarrow T/J\longrightarrow1
 \tag{B.15.10}
\]
is defined, and \(\ker(T\to\mathcal E)=\ker\phi\).
It reconstructs exactly the data consisting of an extension by
\(P\) and a map \(T\to\mathcal E\) onto the quotient, with
\(J\) the inverse image of \(P\). For a nonempty family of such
extensions \(\mathcal E_i\) with identified kernel \(P\) and
quotient \(\gamma\), a compatible transitive system of
isomorphisms \(\mathcal E_i\simeq\mathcal E_j\), fixing both,
is equivalent to a closed subgroup
\[
 \mathcal E\subset\prod_{i,\gamma}\mathcal E_i,
 \qquad \mathcal E\cap P^I=\Delta P,
 \qquad \mathcal E\twoheadrightarrow\gamma.
 \tag{B.15.11}
\]
Such a subgroup is additional data; it need not exist merely
because the quotients and the induced actions agree.
\end{statement}

\begin{proof}
For an automorphism \(u\) of \(C\), a quotient automorphism of
\(C/N\) exists if and only if \(u(N)=N\). It is then unique,
and its generator images and every numerical coordinate are
prescribed. Evaluation on \(N\), both for \(u\) and \(u^{-1}\),
expresses this equality as closed conditions. Hence \(D_q\) is a
closed subgroup. Inner automorphisms preserve \(N\), and \(\tau\)
does so by the compatibility assumption. This proves the assertions
about \(j(L_0)\) and \(t\). Descent is continuous: an open normal
subgroup of \(C'\) pulls back to an open normal subgroup of \(C\),
so all its finite quotient evaluations are continuous. Compactness
gives closedness of \(H\).

Suppose a source effective numerical tuple maps to an effective
target tuple, represented by \(u'\). Then \(qu\) and \(u'q\)
agree on \(\rho,\sigma\), hence on the dense subgroup they
generate, and therefore everywhere. Their kernels give
\(u^{-1}(N)=N\). Conversely kernel stability gives the quotient
automorphism and hence effectiveness of the target tuple. This
proves the three equivalences, without assuming effectiveness of
all numerical solutions. If \(u(L_0)=cL_0c^{-1}\), replace \(u\)
by \(\operatorname{Int}(c^{-1})u\). Both have the same
kernel-stability property since \(N\) is normal. This proves
the version b') of C3:570.

The map \(q:L_0\to L'_0\) is onto: its image is compact and
contains the generator of \(L'_0\). It follows that
\(q_D(j(L_0))=j'(L'_0)\). The compatible map
\(C\rtimes D_q\to C'\rtimes H\) carries one antidiagonal onto
the other, and therefore induces the continuous epimorphism of
the \(G\)'s. It is the unique map with the two specified
compatibilities, since \(C\) and \(D_q\) generate its source.
If a class \(d\in\gamma_{D_q}\) maps to the identity, then
\(q_D(d)=j'(\ell')\). Lifting \(\ell'\) to \(L_0\) and
subtracting its \(j\)-image gives a representative in \(A_q\).
The ambiguity is exactly \(j(L_N)\). This proves the first
formula in (B.15.2).

Similarly any element of \(K_G\), represented in
\(C\rtimes D_q\), can be multiplied by a lifted antidiagonal
element so that its two coordinates lie in \(N\) and \(A_q\).
The intersection of \(N\rtimes A_q\) with the original
antidiagonal is precisely the subgroup indexed by \(L_N\).
This proves the second formula and the exact sequence in
(B.15.2). The assumption \(K=K'=1\) ensures that both copies
of the base group remain embedded; it is not silently discarded.

The signed normalizer map is \((g,e)\mapsto(q_G(g),e)\).
In its kernel the sign is \(+1\), and its group coordinate
centralizes \(x\); this proves (B.15.3). All lifts of \(g'\)
have the form \(kg\), \(k\in K_G\). Such a lift has the
prescribed action on \(x\) exactly when
\(k(gxg^{-1})k^{-1}=x^e\), proving (B.15.4).
Once one lift is chosen, division on the right identifies all
others with \(C_{K_G}(x)\). Kernels commute with these explicitly
defined fibre products, giving (B.15.5). Their images still require
the conjugator-lifting conditions; the proof gives no spurious
surjectivity assertion.

For (B.15.6), the inverse image of \(G_0=CD_0\) in
\(C\rtimes D_0\) has kernel exactly
\(\{(\ell^{-1},j(\ell)):j(\ell)\in D_0\}\).
The quotient by \(C\) is consequently
\(D_0/(D_0\cap j(L_0))\), as asserted. Restricting the
conjugacy arguments in B.13 proves the corresponding normalizer
and fibre-product formulas; the condition \(t\in D_0\) supplies
the same sign splittings. The group \(D_0j(L_0)\) gives the same
\(G_0\) but a different coefficient group, which explains why
these two conventions must not be interchanged without changing
the first kernel. The finite check below distinguishes them.

The effective parameter names of C3:565--569, (70)--(87), also
have an explicit dictionary. With \(P_D^{\mathrm{eff}}\) as
in B.13, its homeomorphism with \(S_0\Gamma_D\) gives
\[
 \Gamma_D\simeq L_0\backslash P_D^{\mathrm{eff}},\qquad
 S_0\mathcal M_D\simeq P_D^{\mathrm{eff}}\times C,
 \qquad f=dU^{-1},\quad U=f^{-1}d.
 \tag{B.15.6a}
\]
These are homeomorphisms of parameter spaces, with the transported
group laws, not an assertion of a direct-product group law in the
second expression. The first follows because multiplication by
\((j(\ell),1,1)\) sends \((\alpha,\beta)\) to
\((\ell\alpha,\ell\beta)\). In the second, equality of the
quotient coordinates puts \(f\) in \(C\), and the displayed
inverse recovers \(U\). The right action
\((\alpha,\beta)\mapsto(\alpha a_0,\beta b_0)\), for
\((a_0,b_0)\in C_\rho\times C_\sigma\), corresponds to
\[
 (d,a,b)\longmapsto
 \bigl(d,(t')^{v_r}a_0^{-1}(t')^{-v_r}a,
          (t')^{v_s}b_0^{-1}(t')^{-v_s}b\bigr).
 \tag{B.15.6b}
\]
This is obtained by inverting (B.13.11). Thus the operator written
\(\tau\) in C3:569, (83), must act here by \(t'\)-conjugation;
literal \(t\)-conjugation need not preserve \(C_\rho\).

We now prove (B.15.8) by identities over
\(\Z[a,\mu,\mu^{-1}]\). Put
\[
 U(z)=\begin{pmatrix}1&z\\0&1\end{pmatrix},\quad
 R=q(u(\rho)),\quad S=q(u(\sigma)).
\]
The conjugacy conditions give \(\operatorname{tr}R=1\),
\(\operatorname{tr}S=0\), and determinants one. In particular
\(S^2=-I\). Since \(e_0=\sigma\rho=U(-1)\), the multiplier
condition gives \(SR=U(-\mu)\). Write
\(S=\left(\begin{smallmatrix}a&b\\c&-a\end{smallmatrix}\right)\).
Then
\[
 R=-S U(-\mu),\qquad
 1=\operatorname{tr}R=c\mu,
 \qquad b=-\mu(1+a^2).
 \tag{B.15.12}
\]
Consequently
\[
 S=\begin{pmatrix}a&-\mu(1+a^2)\\\mu^{-1}&-a\end{pmatrix}
    =g(a,\mu)\sigma g(a,\mu)^{-1},\quad
 R=g(a,\mu)\rho g(a,\mu)^{-1}.
 \tag{B.15.13}
\]
These identities hold also in residue characteristics two and
three. Conjugation by \(g(a,\mu)\in\GL_2(\widehat\Z)\)
is a continuous automorphism of \(\SL_2(\widehat\Z)\).
The two homomorphisms in (B.15.8) agree on the topological
generators, so agree everywhere. Taking kernels proves
\(u(N)=N\). Finally the map
\(\widehat\Z\to L'_0\), \(z\mapsto U(-z)\), is an
isomorphism, and factors through \(L_0\). The surjection
\(\widehat\Z\to L_0\) is thus injective as well. This proves
the asserted kernels. Also \(U(-z)\) commutes with \(\sigma\)
only when \(z=0\), so both required intersections of \(L_0\)
with the centre are trivial in this example.
The surjectivity of (B.15.7) follows from
elementary generation modulo every integer and the Chinese
remainder theorem, as in C.3.

For the last reconstruction, the map
\(j\mapsto(\phi(j)^{-1},j)\) is an injective homomorphism.
The first identity in (B.15.9) proves the homomorphism assertion
and shows that its image centralizes \(P\). The second identity
shows that its image is normalized by \(T\). It is therefore
closed and normal in \(P\rtimes T\). It intersects the first
factor trivially, and intersects the second factor in
\(\ker\phi\); these observations give every assertion in
(B.15.10). Conversely, from a map \(T\to\mathcal E\) onto
\(\mathcal E/P\), conjugation gives \(\theta\), the inverse
image of \(P\) gives \(J\), and restriction gives \(\phi\).
The map \(P\rtimes T\to\mathcal E\) is onto and has exactly
the displayed antidiagonal kernel. The constructions are inverse,
including on compatible morphisms. In particular the notation
\(S_0G=N.\mathcal G\) in C3:692, (4), must mean the specified
semidirect product, except when its action is trivial.

To prove (B.15.11), each coordinate projection
\(\mathcal E\to\mathcal E_i\) is injective: its kernel lies
in \(\Delta P\) and has one coordinate equal to the identity.
It is onto: lift the quotient coordinate of an arbitrary
\(g_i\in\mathcal E_i\) to \(\mathcal E\), then correct its
\(i\)-th coordinate by an element of \(\Delta P\).
Compactness makes each continuous bijection an isomorphism.
Their composites give the desired transitive system. Conversely,
the graph of a transitive system gives such a closed subgroup.
This proves the exact correction made in C3:695 to the tentative
amalgam idea of C3:694.

All the following checks were executed by
\texttt{code/AB\_functoriality.py}; complete output and elementwise
data are in its \texttt{.out} and \texttt{.json} files. Take
\(C=S_4\), \(\rho=(123)\), \(\sigma=(34)\), and let
\(q:S_4\to S_3\) act on the three partitions
\(12|34,13|24,14|23\). Thus
\(\rho'=(132)\), \(\sigma'=(23)\),
\(N=\{1,(12)(34),(13)(24),(14)(23)\}\), and
\(L_N=\{1,(14)(23)\}\). Cayley-graph propagation enumerates
every automorphism and verifies kernel stability. We retain
\(U(12)\), rather than count the infinite numerical group
\(\widehat\Z^\times\), and build every quotient and
fibre product. Each entry below gives source, target, and kernel
orders; every image is the full displayed target.
\[
\begin{array}{c|rrr|rrr}
 &\multicolumn{3}{c|}{D=B}&\multicolumn{3}{c}{D=Z}\\
 \text{object}&|X|&|X'|&|\ker|&|X|&|X'|&|\ker|\\\hline
 D&16&8&2&64&32&2\\
 \gamma&4&4&1&16&16&1\\
 G&96&24&4&384&96&4\\
 N_\rho^*&12&12&1&48&48&1\\
 N_\sigma^*&16&8&2&64&32&2\\
 \Gamma&48&24&2&192&96&2\\
 \mathcal M&1152&144&8&4608&576&8\\
 S_0\Gamma&192&48&4&768&192&4\\
 S_0\mathcal M&4608&288&16&18432&1152&16
\end{array}
\tag{B.15.14}
\]
Here \(A_q=j(L_N)\), \(K_\gamma=1\), and \(K_G=N\), exactly
as (B.15.2) predicts. The program checks all multiplication pairs
for the \(D,\gamma,G\) maps and every image and fibre of all
the listed groups. For \(D_0=\langle t\rangle\) it obtains
\(|D_0|=2\), \(L(D_0)=1\), \(|G_0|=48\), and verifies that
\(t'\notin D\) and that conjugation by \(t'\) does not preserve
\(L_0\). Thus the margin in C3:571 cannot literally require
this \(t'\) in the coefficient group.

The failure of normalizer surjectivity in (B.15.4) is real. Take
\(q:\SL_2(\F_3)\to\SL_2(\F_3)/\{\pm I\}\simeq A_4\)
with the matrices in (B.15.7). The action on
\(\mathbf P^1(\F_3)\) gives the displayed isomorphism: its
kernel is \(\{\pm I\}\), and its generator images are even
permutations, giving all twelve even permutations.
An automorphism acting trivially on this quotient has the form
\(g\mapsto g\chi(g)\), with \(\chi:\SL_2(\F_3)\to C_2\)
a homomorphism. The relations
\(\rho^3=\sigma^2\) and \((\sigma\rho)^3=1\) force
\(\chi(\rho)=\chi(\sigma)=1\), so \(A_q=1\).
Here \(L_N=1\), whence \(K_G=\{\pm I\}\). The element
\[
 j=\begin{pmatrix}1&1\\1&2\end{pmatrix}
 \quad\text{satisfies}\quad j\sigma j^{-1}=-\sigma.
\]
Its image centralizes \(\sigma'\) and has trivial quotient
coordinate, so belongs to the target \(N_{\sigma'}^*\) with
sign \(+1\). Its only possible lifts are \(j,-j\), both of
which fail to centralize \(\sigma\). The program verifies this
obstruction and all candidate lifts. With independent \(U(12)\)
numerics it obtains \(|D|=|H|=24\),
\(|N_\sigma^*|=|N_{\sigma'}^*|=32\), and image order \(16\).
Thus a surjective map of the enlarged groups genuinely can induce
a nonsurjective map of their starred normalizers.

The symbolic calculation verifies (B.15.12)--(B.15.13), and
exhaustive enumeration of determinant-one trace pairs with
\(SR=U(-\mu)\) gives respectively \(2,6,8,20,32,54\)
solutions over \(\Z/n\), for \(n=2,3,4,5,8,9\).
In each case there are exactly \(n\) pairs for each unit
\(\mu\), parameterized by \(a\in\Z/n\). These are trace-pair
counts, not an assertion that every independently signed
conjugacy condition is automatic over every finite ring.

For the reconstruction and gluing test, take \(P=\gamma=C_2\)
with trivial action. The groups \(C_4\) and \(V_4\) are two
extensions with that kernel and quotient; their different element
orders prove that they cannot be isomorphic. The program reconstructs
both by (B.15.10). It then enumerates every subgroup of order four
in each relevant fibre product of order eight. For two copies of
\(C_4\), precisely two subgroups satisfy (B.15.11), the graphs
of its two automorphisms. For \(C_4\) and \(V_4\), there are
none. This is a complete finite counterexample to reconstruction
from only the common quotient and action, and verifies the
additional gluing data required in C3:695--696. Passing to constant
finite group schemes makes these particular kernel and quotient
statements exact in fppf sheaves over \(\Z\).
\end{proof}

\paragraph{The congruence comparison left by this theorem.}
For (B.15.7), write \(H_q=q_D(D)\subset D'\) and
\(A_q=\ker(D\to D')\). The theorem constructs the maps and
determines their kernels in terms of these two well-defined groups;
it does not determine the full effective image \(H_q\) or the
group \(A_q\). Its elements have numerical coordinates
\((\mu,e_r,e_s)=(1,1,1)\), and satisfy \(qu=q\); the coordinates
have not been forgotten before taking the kernel. In particular a finite target tuple must lift
compatibly to an automorphism of the full completed modular group,
not just satisfy its matrix identities. B.8.5, the named comparison
of B.16, the restricted-conjugator comparisons of B.18, and the
three arithmetic comparison problems in D.7 are separate assertions.
The finite map (B.15.14) is a proved nontrivial descent and
conjugator-lifting case, not a substitute for any of them.

Here is one completely specified lifting instance not supplied by
our present argument. Let \(\mu\in\widehat\Z^\times\) have
\(5\)-adic component \(2\) and every other component \(1\),
and put
\[
 d'_5=\bigl(\operatorname{Int}(\operatorname{diag}(\mu,1)),
                     \mu,1,1\bigr)\in Z'.
 \tag{B.15.15}
\]
It belongs also to \(B'\). Indeed, both prescribed signs are
\(+1\). At \(5\), choose \(i\in\Z_5\) with \(i^2=-1\),
\(i\equiv2\pmod5\); lifting the simple root modulo successive
powers of five constructs it. The centralizers contain
\[
 h_r=(i-1)I+2\rho,\quad \det h_r=2,
 \qquad h_s=I+\sigma,\quad\det h_s=2.
\]
Thus \(\operatorname{diag}(2,1)h_x^{-1}\in\SL_2(\Z_5)\)
conjugates \(x\) to the required image for \(x=\rho,\sigma\).
At every other prime take the identity conjugator. The loop condition
is immediate from conjugation of \(U(-1)\). The code verifies the
two resulting \(\SL_2(\F_5)\) conjugators, namely
\(\left(\begin{smallmatrix}3&3\\1&3\end{smallmatrix}\right)\)
and \(\left(\begin{smallmatrix}1&1\\2&3\end{smallmatrix}\right)\).
The precise assertion left by this example is
\[
 \exists u\in\Aut\bigl(\widehat{\SL_2(\Z)}\bigr):\quad
 (u,\mu,1,1)\in Z,
 \qquad qu=\operatorname{Int}(\operatorname{diag}(\mu,1))q.
 \tag{B.15.16}
\]
Equivalently, \(d'_5\in H_q\). The separate conjugacy and norm
conditions just proved do not establish this compatible profinite
realization. We record (B.15.16) as an exact unresolved instance of
our argument, not as a claim that it is open in the literature or
the least unit for any order on \(\widehat\Z^\times\).
